\begin{figure*}[!h]
\centering
\setlength{\tabcolsep}{1.5pt}

    \newcommand{\resultfigwidth}{.193\textwidth}
	\begin{tabular}{cc|cccc}
		& Ground Truth & 2DGS* & BBSplat & GStex & Ours \\
        \raisebox{0.25\height}{\rotatebox{90}{\small Mip-Nerf360}} &
        \spyimage{red}{\resultfigwidth}{images/table1/m360_gt.jpg}{0.4, 1.5}{1.1, 0.6} &
        \spyimage{red}{\resultfigwidth}{images/table1/m360_2dgs_no_regul.jpg}{0.4, 1.5}{1.1, 0.6} &
        \spyimage{red}{\resultfigwidth}{images/table1/m360_bbsplat_baseline.jpg}{0.4, 1.5}{1.1, 0.6} &
        \spyimage{red}{\resultfigwidth}{images/table1/m360_gstex_output_1e7.jpg}{0.4, 1.5}{1.1, 0.6} &
        \spyimage{red}{\resultfigwidth}{images/table1/m360_texture-proj_final_new_downscale_threshold_0.02.jpg}{0.4, 1.5}{1.1, 0.6} \\
		
        \raisebox{0.18\height}{\rotatebox{90}{\small Deep Blending}} &
        \spyimage{red}{\resultfigwidth}{images/table1/db_gt.jpg}{3, 0.3}{3.3, 1.3} &
        \spyimage{red}{\resultfigwidth}{images/table1/db_2dgs_no_regul.jpg}{3, 0.3}{3.3, 1.3} &
        \spyimage{red}{\resultfigwidth}{images/table1/db_bbsplat_baseline.jpg}{3, 0.3}{3.3, 1.3} &
        \spyimage{red}{\resultfigwidth}{images/table1/db_gstex_output_1e7.jpg}{3, 0.3}{3.3, 1.3} &
        \spyimage{red}{\resultfigwidth}{images/table1/db_texture-proj_final_new_downscale_threshold_0.02.jpg}{3, 0.3}{3.3, 1.3} \\

        \rotatebox{90}{\small Tanks\&Temples} &
        \spyimage{red}{\resultfigwidth}{images/table1/tnt_gt.jpg}{3.1, 1.65}{3.1, 0.6} &
        \spyimage{red}{\resultfigwidth}{images/table1/tnt_2dgs_no_regul.jpg}{3.1, 1.65}{3.1, 0.6} &
        \spyimage{red}{\resultfigwidth}{images/table1/tnt_bbsplat_baseline.jpg}{3.1, 1.65}{3.1, 0.6} &
        \spyimage{red}{\resultfigwidth}{images/table1/tnt_gstex_output_1e7.jpg}{3.1, 1.65}{3.1, 0.6} &
        \spyimage{red}{\resultfigwidth}{images/table1/tnt_texture-proj_final_new_downscale_threshold_0.02.jpg}{3.1, 1.65}{3.1, 0.6} \\
	\end{tabular}
\caption{
\label{fig:default-eval}
We provide renderings from one scene per dataset for every method,
trained with default settings.
Our method is able to reconstruct high frequency details on images,
even while using fewer parameters compared to the other methods.
}
\end{figure*}



\section{Implementation}
We have implemented our method using the 3DGS codebase, introducing 2D primitives as in 2DGS. 
%\PP{we used 3dgs and then added bits of 2dgs}.
Unlike 2DGS,
we do not use the $\mathbb{R}^2 \rightarrow \mathbb{R}^2$ transformation to find the intersection point using three non-parallel planes.
Instead, we use Eq.~\ref{eq:intersection_point} directly in camera space.
We did not observe any instability issues.
\revision{rev6315 runtime performance}{
Since our method requires per-ray querying of texture maps,
 both the forward and backward passes incur additional overhead,
that is not compensated by the smaller number of primitives rendered.
In general, compared to 2DGS,
 training time is 1.5-2 times longer, and rendering is 25\% slower.
}

As the texture grids can have different resolutions,
we created custom ``jagged'' tensor data structure to be used for their storage.
The texture resolution for each primitive gets (de-)allocated dynamically,
so that it covers the extent of the primitive,
that is $\pm$3 standard deviations.
This dynamic memory management happens every 100 iterations.
A hard limit of 256 texture resolution is enforced to avoid using too much memory.
When a primitive grows outside of its allocated texture grid,
either because memory has not been allocated yet or it has surpassed the hard limit,
zero-padding is used. Since we are storing an offset, zero-padding reduces to the DC color.

We observed that limiting the $t_2p_r$ to be greater or equal to 2 did not have a detrimental effect on the visual quality,
as subtexel scaled details can be retrieved
thanks to alpha blending and the overlap of our primitives.
%So,
We thus allow upscaling to happen only for primitives with a $t_2p_r$ greater or equal to 2.
% The error calculation,
% adaptive texel size determination and resolution-aware primitive management strategies run every 250 iterations,
% until 25k iterations.
% The threshold $\tau_{\mathrm{tr}}$ starts at $64$ and is progressively reduced to $32$ within 7000 iterations.

% The texture grids are initialized with a value of $0$ on every channel
% and start optimizing after 500 iterations.
% To avoid running out of memory and to avoid early overfitting,
% at 500 iterations,
% we set $t_2p_r$ for each primivite so that their smallest axis is 8 texels wide,
% in practice using the closest power of two value.

% We implement downscaling and upscaling using pytorch's interpolate function,
% with a scale factor of 2,
% which translates to reductions or increases of the number of texels by a factor of 4,
% respectively.
% Upscale is done using nearest neighbour interpolation,
% so that we avoid changing the appearance of the texture,
% that could disrupt that optimizer.



\revision{rev4160 memory/storage}
{
\textbf{Storage and Memory Considerations.}
The size of our representation is tied to the number of parameters used,
with each parameter saved as a 4-byte float,
as in previous Gaussian Splatting implementations.
However,
%as our representation constitutes an extension to the original one,
%we expect 
most Gaussian Splatting compression techniques (see \cite{3dgszip}) are 
applicable to our approach either directly or with minor modifications.
For our newly introduced texture maps,
our choice to use a sigmoid activation to limit the effective range of the texel values allows for a very efficient application of K-means clustering compression. Please see Tab.~\ref{tab:per-scene} for complete statistics.
}